\chapter{Elementary Definitions and Conventions}

\section{Rings and Ideals}

\begin{definition}{Rings}\\
A \textit{ring} is an abelian additive group with a multiplicative operation $(a,b) \mapsto ab$ satisfying the associativity and distributivity. Moreover, we say that a ring $R$ is
\begin{itemize}
    \item \textit{commutative} if for any $a,b \in R$, $ab=ba$;
    \item \textit{unital} if there exists an multiplicative identity $1_R$. Furthermore, an element $u$ in a unital ring $R$ is called a \textit{unit} (or an \textit{invertible element}) if there exists $v \in R$, s.t.: $uv=1_R$, which is denoted as $v=u^{-1}$ and called the \textit{inverse} of $u$. We denote by $U(R)$ the set of all units in $R$ and then we get the definition of \textit{field}: a commutative unital ring $F$ is a \textit{field} if $F \setminus \{0_F\}=U(F)$.
\end{itemize}
\end{definition}

\noindent\textit{Convention:} In this book, every ring is treated as a commutative unital ring. Furthermore, in this book, unless otherwise specified, the letter "$R$" represents a (commutative unital) ring.

\begin{definition}{Ideals}
\begin{itemize}
    \item A subset $I \subseteq R$ of $R$ is called an \textit{ideal} of $R$, denoted as $I \triangleleft R$, if $I$ is additive subgroup of $R$ and if for any $r\in R$, any $x\in I$, we have $rx \in I$.

    \item An ideal $I \triangleleft R$ is said to be \textit{generated by a subset} $S \subseteq R$ if for any $x \in I$, we can write $x$ in the form
    $$x= \sum_{i=1}^n r_is_i$$
    where $r_i \in R, s_i \in S$. Then we denote $I$ as $(S)$.
    
    \noindent\textit{Convention:} We write as $0$ the ideal generated by the empty set $\emptyset$.

    \item An ideal is \textit{principal} if it can be generated by one single element.
    \item An ideal $I \triangleleft R$ is \textit{prime} if $I$ is proper, i.e.: $I \neq R$, and if $xy \in I$ implies $x \in I$ or $y \in I$.
    \item An ideal is \textit{maximal} if it is proper and if it is NOT contained in any other proper ideal.
\end{itemize}
\end{definition}

\begin{proposition}{Prime Ideal-Maximal Ideal Proposition}
\begin{enumerate}
    \item $I \triangleleft R$ is prime $\Leftrightarrow$ $\forall$ ideals $J,K \triangleleft R$ with $J,K \subseteq I$, we have $J \subseteq I$ or $K \subseteq I$ $\Leftrightarrow$ $R/I$ is an ID.
    \item $I \triangleleft R$ is maximal $\Leftrightarrow$ $R/I$ is a field.
\end{enumerate}
\end{proposition}

\begin{proof}
    Omitted.
\end{proof}

\begin{definition}{Domain}\\
$R$ is called a \textit{domain} if $0= (\emptyset)$ is prime.
\end{definition}

\begin{definition}{Prime Elements}\\
An element $p \in R$ if \textit{prime} if $(p)$ is a prime ideal, or equivalently, if $p \notin U(R)$ and if $p \mid xy$ $\Rightarrow$ $p\mid x$ or $p\mid y$.
\end{definition}

\begin{definition}{Ring Homomorphisms}\\
A map $\varphi: R \to S$ from a ring $R$ to a ring $S$ is called a \textit{ring homomorphism} if it preserves addition and multiplication. Furthermore, if $\varphi(1_R)= 1_S$, then we say $\varphi$ is \textit{unital}.
\end{definition}

\noindent\textit{Convention:} In this book, we will omit the adjectives "unital" and "ring", i.e.: "homomorphism" means "unital ring homomorphism".

\begin{definition}
    Let $r,s,r_1,\cdots,r_n \in R$ be elements of $R$.
    \begin{itemize}
        \item $r$ is an \textit{idempotent} if $r^2=r$.
        \item $r$ and $s$ are \textit{orthogonal} if $rs=0$.
        \item $\{r_1,\cdots,r_n\}$ is a \textit{complete set of orthogonal idempotents} in $R$ if every $r_i$ is an idempotent in $R$, if $\forall i\neq j \in \{ 1,\cdots, n\}$, $r_ir_j =0$, i.e.: they are orthogonal, and if $\sum_{i=1}^n r_i=1_R$.
    \end{itemize}
\end{definition}

\noindent\textbf{Example:} $\{(1_{R_1}, 0,\cdots, 0) , (0,1_{R_2},0,\cdots, 0) ,\cdots, (0,\cdots, 0,1_{R_n}) \} \subseteq R_1 \times \cdots \times R_n$ is a \textit{complete set of orthogonal idempotents} in $R_1 \times \cdots \times R_n$.

\begin{proposition}{Complete Set of Orthogonal Idempotents-Decomposition Proposition}\\
If $\{e_1 \cdots, e_n \}$ is a complete set of orthogonal idempotents in $R$, then
$$R \cong R e_1 \times \cdots \times Re_n$$
is a direct product decomposition.
\end{proposition}

\begin{proof}
    We only deal with the case $n=2$. Consider the map $\varphi: R \to R e_1 \times Re_2$ given by $\varphi(r):= (re_1,re_2)$.
    \begin{itemize}
        \item \textit{$\varphi$ is a homomorphism:} For any $r,s \in R$, let $r=r_1e_1 +r_2e_2$ and $s=s_1e_1+ s_2e_2$.
        \begin{itemize}
            \item $\varphi(r+s) := (r_1e_1+s_1e_1 ,r_2e_2 + s_2e_2) = (r_1e_1,,r_2e_2) +(s_1e_1, s_2e_2) := \varphi(r)+\varphi(s)$.
            \item $\varphi (rs) := ((rs)e_1,(rs)e_2) = ((re_1)(se_1) ,(re_2)(se_2)) = (re_1,se_1) (re_2,se_2) := \varphi(r)\varphi(s)$.
        \end{itemize}

        \item \textit{$\varphi$ is injective:} $\ker \varphi :=\{ r\in R \mid re_1=0,re_2=0 \} =\{r\in R\mid r_1=0,r_2=0\}= \{0\}$.
        \item \textit{$\varphi$ is surjective:} $\forall (x,y) \in Re_1 \times Re_2$, $\exists r_1,r_2 \in R$, s.t.: $x=r_1e_1,y=r_2e_2$. Then $\varphi(r_1e_1+r_2e_2) := ((r_1e_1+r_2e_2)e_1,(r_1e_1+r_2e_2)e_2) = (r_1e_1,r_2e_2) = (x,y)$.
    \end{itemize}
\end{proof}

\noindent\textbf{Remark:} Conversely, if $R \cong R_1 \times \cdots \times R_n$ can be decomposed as a direct product of finitely many (nontrivial) subrings, then $\{ e_i:= (0,\cdots,0,\underbrace{1}_{ith\; position},0,\cdots,0)\}_{i=1}^n$ is a complete set of orthogonal idenpotents.

\begin{definition}{Commutative Algebra}
\begin{itemize}
    \item A \textit{commutative algebra over $R$} (or \textit{commutative $R$-algebra}) is a ring $S$ together with a homomorphism $\alpha: R \to S$, namely $(S,\alpha)$.

    \noindent\textit{e.g.:} \begin{itemize}
        \item Any ring $S$ together with the homomorphism $\alpha: n \mapsto \begin{cases}
            \overbrace{1+ \cdots +1}^{n-times} &,if\;n>0\\
            0 &,if\; n=0\\
            \underbrace{(-1) + \cdots + (-1)}_{(-n)-times} &,if \;n<0
        \end{cases}$ is a commutative $\mathbb{Z}$-algebra. Furthermore, $(S,\alpha)$ is the \textit{unique} commutative $\mathbb{Z}$-algebra.

        \item $(R[x_1,\cdots,x_n] , \substack{R \;\to\; R[x_1,\cdots,x_n]\\
        r \;\mapsto\; r+ 0x_1 +\cdots+ 0x_n})$ is a commutative $R$-algebra.
    \end{itemize}

    \item We usually omit the homomorphism $\alpha$ and write $rs:= \alpha(r) s$ for $r\in R,s\in S$.
    \item A \textit{subalgebra} of $S$ is a subring $S'$ of $S$ that contains $\alpha(R)$.

    \noindent\textit{e.g.:} $R[x_1^2,x_2^3]$ is a $R$-subalgebra of $R[x_1,x_2]$.
    
    \item A \textit{homomorphism of $R$-algebra} $\varphi:S \to T$ is a homomorphism of rings such that $\varphi(rs) = r\varphi(s)$ $\forall r\in R,s\in S$. More precisely, $\varphi: (S,\alpha_S) \to (T,\alpha_T)$ is a homomorphism of rings from $S$ to $T$ such that $\varphi (\alpha_S (r) s) =\alpha_T (r)\varphi(s)$ $\forall r\in R,s\in S$.

    \noindent\textit{e.g.:} Given $r\in R$, $\varphi : \substack{R[x]\;\to\; R\\
    f \;\mapsto\; f(r)}$ is a homomorphism of $R$-algebras.
\end{itemize}
\end{definition}

\begin{definition}{Polynomial Rings}\\
Let $R[x_1,\cdots,x_n]$ be a \textit{polynomial ring over $R$}.
\begin{itemize}
    \item The elements in $R$ are referred to as \textit{scalars}.
    \item A \textit{monomial} is a product of variables, whose \textit{degree} is defined as the number of these factors (counting repeats), e.g.: $\deg (x_1^2x_2^3) = \deg (x_1x_1x_2x_2x_2) =5$. By \textit{convention}, we regard $1$ as the empty product, i.e.: $1$ is the unique monomial of degree $0$.

    \item A \textit{term} is a scalar timing a monomial, i.e.: of the form:
    $$r \cdot x_1 ^{a_1} \cdots x_n^{a_n} \;\;where\;r\in R, a_i \in \mathbb{N} \cup \{0\}$$
    \item $f\in R[x_1,\cdots ,x_n]$ is said to be \textit{homogeneous} if $f=0$ or if the monomials in the terms of $f$ all have the same degree. We also use the word "\textit{form}" to mean "homogeneous polynomial".
\end{itemize}
\end{definition}

\section{Unique Factorization}

\noindent\textit{Convention:} In this book, we say that $R$ is \textit{factorial} if it is a UFD.

\begin{theorem}
    PID $\Rightarrow$ FD $\Rightarrow$ UFD.
\end{theorem}

\begin{proof}
    Outline: PID $\Rightarrow$ ACCP $\Rightarrow$ FD.
\end{proof}

